\chapter{Lecture 15: Method Overriding, Virtual Methods and Polymorphism}
\begin{center}
\textit{March 1, 2018}
\end{center}

\section{Overriding}
We want to decide whether a book is heavy: a Book is heavy if it has more than
200 pages, a Text if more than 500, a Comic if more than 30.

\begin{lstlisting}[style=cpp]
class Book {
    ...
    int numPages;
  public:
    int getNumPages() const { return numPages; }
    bool isHeavy() const { return numPages > 200; }
};

class Comic : public Book {
    ...
  public:
    bool isHeavy() const { return getNumPages() > 30; }   // overriding
};
// Text is similar, with 500.

Book b{"title", "author", 50};
cout << b.isHeavy();     // Book::isHeavy: false

Comic c{"title", "author", 40, "hero"};
cout << c.isHeavy();     // Comic::isHeavy: true

Book b1 = Comic{"title", "author", 40, "hero"};   // compiles!
cout << b1.isHeavy();    // Book::isHeavy runs: false
\end{lstlisting}

\begin{definition}
\textbf{Method overriding}: a subclass replacing the implementation of an
inherited method. (Not to be confused with \emph{overloading}: same name,
different parameter lists.)
\end{definition}

\section{Object Slicing}
In \code{Book b1 = Comic\{...\};}, \code{b1} is a Book: there is only enough
space for a Book's fields. The Comic is copied using Book's copy ctor, and the
\code{hero} field is chopped off. This is called \textbf{object slicing}
(another professor calls it \emph{object coercion}). So \code{b1} really is a
Book, and \code{Book::isHeavy} runs.
\begin{center}
\begin{tikzpicture}[>=Stealth,font=\small,f/.style={draw,minimum width=2.1cm,minimum height=0.55cm}]
  \node[font=\scriptsize] at (0,0.55) {\code{Comic\{...\}}};
  \node[f,fill=blue!6] (t) at (0,0) {\code{title}};
  \node[f,fill=blue!6,below=0pt of t] (a) {\code{author}};
  \node[f,fill=blue!6,below=0pt of a] (n) {\code{numPages}};
  \node[f,fill=green!10,below=0pt of n] (h) {\code{hero}};
  \node[font=\scriptsize] at (4.5,0.55) {\code{Book b1}};
  \node[f,fill=blue!6] (t2) at (4.5,0) {\code{title}};
  \node[f,fill=blue!6,below=0pt of t2] (a2) {\code{author}};
  \node[f,fill=blue!6,below=0pt of a2] (n2) {\code{numPages}};
  \draw[->] (a.east) -- node[above,font=\scriptsize]{Book's copy ctor} (a2.west);
  \draw[red!70!black,thick] ([xshift=0.3cm]h.east) -- ++(0.4,0.25) ([xshift=0.3cm,yshift=0.25cm]h.east) -- ++(0.4,-0.25);
  \node[font=\scriptsize,red!70!black,anchor=west] at ([xshift=0.8cm]h.east) {sliced off};
\end{tikzpicture}
\end{center}

Slicing does not occur if a subclass object is pointed to by a superclass
pointer (or reference):
\begin{lstlisting}[style=cpp]
Comic c{"title", "author", 40, "hero"};
Comic *cp{&c};
cout << cp->isHeavy();   // Comic::isHeavy: true
Book *bp{&c};            // fine: a Comic is-a Book
cout << bp->isHeavy();   // Book::isHeavy runs: false!
\end{lstlisting}
The Comic object only acts like a Comic when a Comic pointer points to it. For
non-virtual methods, the compiler looks at the \emph{declared} type of the
pointer and decides at compile time which method to run.

But most of the time this is not what we want:
\begin{lstlisting}[style=cpp]
Book *myCollection[100];
...
for (int i = 0; i < 100; ++i) {
    cout << myCollection[i]->isHeavy() << endl;   // always Book::isHeavy
}
\end{lstlisting}
We want the ``most specific'' \code{isHeavy} to run for each book.

\section{Virtual Methods}
\begin{lstlisting}[style=cpp]
class Book {
    ...
  public:
    virtual bool isHeavy() const { return numPages > 200; }
};

class Comic : public Book {        // same for Text
    ...
  public:
    bool isHeavy() const override { return getNumPages() > 30; }
};
\end{lstlisting}
The keyword \code{override} is optional, but recommended: if you make a mistake
in the signature (e.g.\ forget the \code{const}), the method would silently be a
new, unrelated method; with \code{override} it won't compile.

\begin{lstlisting}[style=cpp]
Comic c{"title", "author", 40, "hero"};
Comic *cp{&c};
Book *bp{&c};
Book &rb{c};
cout << cp->isHeavy();   // Comic::isHeavy
cout << bp->isHeavy();   // Comic::isHeavy
cout << rb.isHeavy();    // Comic::isHeavy
\end{lstlisting}
\begin{center}
\begin{tikzpicture}[>=Stealth,font=\small]
  \node[draw,fill=blue!5,minimum width=1.4cm,minimum height=0.7cm] (bp) at (0,0) {\code{bp}};
  \node[font=\scriptsize,anchor=north] at (bp.south) {declared type \code{Book*}};
  \node[draw,fill=green!10,minimum width=2.2cm,minimum height=0.9cm] (c) at (3.4,0) {Comic object};
  \draw[->] (bp) -- (c);
  \node[draw,rounded corners] (b) at (8,0.7) {\code{Book::isHeavy}};
  \node[draw,rounded corners] (cc) at (8,-0.7) {\code{Comic::isHeavy}};
  \draw[->,gray] (c.east) -- node[above,sloped,font=\scriptsize]{non-virtual} (b.west);
  \draw[->,thick,blue!60!black] (c.east) -- node[below,sloped,font=\scriptsize]{virtual} (cc.west);
\end{tikzpicture}

\small Static dispatch follows the pointer's declared type; dynamic dispatch follows the object.
\end{center}

\begin{definition}
For a \textbf{virtual} method, the method to call is chosen while the program
is running, based on the \emph{run-time type} of the object, not the declared
type of the pointer/reference. This is \textbf{dynamic dispatch}.
\end{definition}
\begin{itemize}
  \item Dynamic dispatch has a (small) run-time cost (see the chapter on
        vtables).
  \item C++'s default is \textbf{static dispatch} (non-virtual): you don't pay
        for what you don't use.
\end{itemize}

The run-time type is genuinely not known at compile time:
\begin{lstlisting}[style=cpp]
Book *bp;
if (userSaysComic) bp = &c; else bp = &t;
cout << bp->isHeavy();   // decided at run time

Book *myBooks[100];      // store Books, Comics and Texts
for (int i = 0; i < 100; ++i) {
    cout << myBooks[i]->isHeavy() << endl;   // the right one for each
}
\end{lstlisting}
\code{myBooks} is called a \emph{polymorphic array}.
\begin{center}
\begin{tikzpicture}[>=Stealth,font=\small,c/.style={draw,minimum width=1cm,minimum height=0.65cm}]
  \foreach \i in {0,...,3} \node[c] (p\i) at (\i,0) {};
  \node[c] at (4,0) {$\cdots$};
  \node[anchor=east] at (-0.6,0) {\code{myBooks}};
  \foreach \i/\t/\col/\x in {0/Book/blue!6/-1.2,1/Text/orange!15/1.2,2/Comic/green!12/3.6,3/Text/orange!15/6.0} {
    \node[draw,fill=\col,rounded corners,minimum width=1.6cm,align=center] (o\i) at (\x,-1.7) {\t\\ \scriptsize\code{\t::isHeavy}};
    \draw[->] (p\i.center) -- (o\i.north);
  }
\end{tikzpicture}
\end{center}

\begin{definition}
\textbf{Polymorphism} (``many forms''): the ability of one variable/container to
accommodate objects of multiple types.
\end{definition}
We have used it all along: in
\begin{lstlisting}[style=cpp]
ostream &operator<<(ostream &out, const Vec &v);
\end{lstlisting}
\code{out} is a polymorphic parameter: we can pass \code{cout}, an
\code{ofstream}, or an \code{ostringstream}.
\begin{lstlisting}[style=cpp]
Vec v{1, 2};
cout << v;
ofstream file{"vec.txt"};
file << v;               // same operator<<
\end{lstlisting}

\section{Destructors Revisited}
Object destruction is the reverse of object construction. When a subclass
object is destroyed:
\begin{enumerate}
  \item the subclass dtor body runs;
  \item fields of the subclass that are objects are destroyed;
  \item the superclass part is destroyed (the superclass dtor is called
        automatically);
  \item the space is released.
\end{enumerate}

\begin{lstlisting}[style=cpp]
// c++/inheritance/example5
class X {
    int *x;
  public:
    X(int n) : x{new int[n]} {}
    ~X() { delete [] x; }
};

class Y : public X {
    int *y;
  public:
    Y(int n, int m) : X{n}, y{new int[m]} {}
    ~Y() { delete [] y; }      // X's dtor is called automatically after
};

Y myY{10, 20};          // fine: ~Y then ~X at end of scope
X *myX = new Y{10, 20}; // fine: a Y is-a X
delete myX;             // calls ~X() only!
\end{lstlisting}
The dtor is non-virtual, so \code{delete myX} uses the declared type: only
\code{\textasciitilde X} runs, and \code{y}'s array is leaked. Fix: make the base
class dtor virtual (which makes all subclass dtors virtual too):
\begin{center}
\begin{tikzpicture}[>=Stealth,font=\small,f/.style={draw,minimum width=1.4cm,minimum height=0.6cm}]
  \node[f,fill=blue!6] (x) at (0,0) {\code{x}};
  \node[f,fill=green!10,below=0pt of x] (y) {\code{y}};
  \node[font=\scriptsize,anchor=south] at (x.north) {\code{*myX} (a Y)};
  \node[draw,fill=blue!6,minimum width=2.6cm] (ax) at (4,0) {\code{int[10]}};
  \node[draw,fill=red!12,draw=red!60!black,minimum width=2.6cm] (ay) at (4,-0.9) {\code{int[20]}};
  \draw[->] (x.east) -- (ax.west); \draw[->] (y.east) -- (ay.west);
  \node[font=\scriptsize,anchor=west] at (5.4,0) {freed by \code{\textasciitilde X}};
  \node[font=\scriptsize,anchor=west,red!60!black] at (5.4,-0.9) {leaked: \code{\textasciitilde Y} never runs};
\end{tikzpicture}
\end{center}
\begin{lstlisting}[style=cpp]
class X {
    ...
  public:
    virtual ~X() { delete [] x; }
};
\end{lstlisting}

\begin{important}[Rule]
Always make the base class dtor virtual, even if its body is empty, if the
class might ever have subclasses.
\end{important}

If a class should never be subclassed, declare it \code{final}:
\begin{lstlisting}[style=cpp]
class Y final : public X { ... };   // Y can't be subclassed
class Z : public Y { ... };         // won't compile
\end{lstlisting}
\begin{quote}
\textit{\{C++: gotta love it.\}} \code{final} and \code{override} are only
keywords in these particular positions; you can still name a variable
\code{override}.
\end{quote}

\section*{Summary}
\begin{itemize}
  \item Overriding replaces an inherited method's implementation.
  \item Assigning a subclass object to a superclass \emph{object} slices it;
        pointers and references don't slice.
  \item Non-virtual methods use static dispatch (declared type); virtual methods
        use dynamic dispatch (run-time type). Mark overrides with
        \code{override}.
  \item Polymorphism: one pointer/reference type, many object types.
  \item Base class dtors must be virtual; use \code{final} to forbid
        subclassing.
\end{itemize}
